\section{Introduction}
\label{sec:introduction}

The safe operation of a nuclear power plant depends on a control-room decision process that cascades across three questions in rapid succession. First, the operator must decide whether plant behaviour has departed from normal operating conditions at all. Second, if a departure has occurred, its type must be identified: a loss-of-coolant accident (LOCA), a steam-generator tube rupture (SGTR), a main steam line break (MSLB), and many other failure modes evolve with distinct physical signatures and demand different emergency operating procedures. Third, the operator must anticipate how the plant will evolve over the next several minutes so that mitigating actions can be selected before safety limits are challenged. Each of these questions maps onto an established class of data-driven problems---anomaly detection, multi-class classification, and time-series forecasting---but in practice they are rarely addressed as a coordinated system.

Most machine-learning studies on nuclear fault diagnosis focus on a single stage of this chain. Purely data-driven classifiers achieve strong in-distribution accuracy on simulator data but can produce predictions that are inconsistent with reactor physics when deployed on out-of-distribution transients~\citep{qi2022nppad,zhang2023pinn,santhosh2019survey}. Surrogate reactor models trained with physics-informed losses can preserve conservation laws but are typically evaluated as stand-alone models rather than as components of a broader diagnostic system~\citep{raissi2019pinn,zhang2023pinn,gong2022rom}. Recent digital-twin review articles explicitly call for integrated monitoring--diagnosis--prognosis architectures~\citep{ayodeji2022dtwin,digitaltwin2025review}, yet few published works deliver concrete implementations that span the full chain.

This paper contributes such an implementation. We build a three-stage deep-learning framework for PWR fault diagnosis on the open-source Nuclear Power Plant Accident Dataset (\nppad)~\citep{qi2022nppad}, which provides 96 sensor channels across 18 operating conditions generated by the PCTRAN PWR simulator. The three stages and their design rationale are summarized as follows:

\begin{enumerate}[leftmargin=*]
    \item \textbf{Semi-supervised anomaly detection.} We train three architectures (autoencoder, LSTM next-step predictor, masked Transformer reconstructor) on normal operating data only, reflecting the practical reality that accident data at a real plant is scarce, simulated, or both. Anomaly scores are thresholded on validation data. This stage acts as a broad screening gate.
    \item \textbf{Multi-class accident classification with root-cause attribution.} We compare a 1D-CNN, a bidirectional LSTM, and a Transformer with a learnable classification token on the 18-class problem, and we apply SHAP and input $\times$ gradient attribution to the best-performing model to recover per-sensor and per-timestep importance scores. Crucially, we interpret these scores against the known physics of each accident type: for well-understood classes such as LOCA, SLB, and SGTR, we check whether the attributions highlight the sensors predicted by the underlying physical mechanism.
    \item \textbf{Physics-informed trajectory prediction.} We train a GRU encoder--decoder surrogate with a composite loss that combines mean-squared trajectory error, a point-kinetics residual, an energy-conservation residual, and physical-bounds penalties. The surrogate predicts ten time steps ahead from a twenty-step context window, allowing us to generate a short-horizon forecast that can be compared against the observed trajectory as a physics-grounded consistency check.
\end{enumerate}

The central novelty of the paper is not any single stage but the way they are composed. Each stage addresses a distinct question, and the output of one stage conditions the input to the next. More importantly, the physics-informed surrogate of stage three provides a mechanism by which the diagnostic outputs of stages one and two can be cross-validated against reactor physics: if the predicted trajectory conditioned on a classifier's decision is physically inconsistent with what actually unfolds, the classification can be flagged as low-confidence. This cross-validation loop is, to our knowledge, absent from prior published work on \nppad{}.

Our contributions are:

\begin{itemize}[leftmargin=*]
    \item A reproducible three-stage fault diagnosis pipeline on \nppad{} covering anomaly detection, 18-class accident classification, and physics-constrained trajectory prediction, with open configurations and training scripts.
    \item A quantitative evaluation showing that three independent anomaly-detection architectures achieve near-perfect screening ($\mathrm{F1} \approx 0.9995$), that a bidirectional LSTM is the strongest classifier with $71.2\%$ accuracy and $0.572$ macro-F1 on the imbalanced 18-class problem, and that embedding point kinetics and energy conservation into the surrogate's loss reduces physics residuals by $68.6\%$ and energy-conservation violations by $39.6\%$ relative to an unconstrained baseline.
    \item A physics-grounded interpretation of SHAP attributions for representative accident classes, establishing that the classifier's decision process is aligned with mechanistic knowledge of the underlying transients.
    \item A pipeline-integration analysis showing how physics-informed predictions can serve as a consistency check on data-driven classifications, together with a discussion of the conditions under which such cross-validation is most informative.
\end{itemize}

The remainder of the paper is organized as follows. Section~\ref{sec:related} reviews related work. Section~\ref{sec:method} describes each of the three stages and their integration. Section~\ref{sec:setup} details the experimental protocol. Section~\ref{sec:results} presents quantitative and qualitative results. Section~\ref{sec:discussion} discusses limitations and generalization to real plants and advanced reactor designs. Section~\ref{sec:conclusion} concludes.
